% Part: first-order-logic
% Chapter: lindstrom
% Section: introduction

\documentclass[../../../include/open-logic-section]{subfiles}

\begin{document}

\olfileid{mod}{lin}{int}
\section{परिचय}

काही आणखी निर्बंध गृहीत धरल्यास संहतता आणि अधोगामी
लोव्हेनहाइम--स्कोलेम प्रमेये ज्या महत्तम तर्कशास्त्रासाठी खरी ठरतात, ते
प्रथम-क्रम तर्कशास्त्र आहे, असे सांगणारे लिंडस्ट्रॉमचे लक्षणचित्रण या
प्रकरणात सिद्ध करणे हे आपले उद्दिष्ट आहे
(\olref[fol][com][com]{thm:compactness} आणि
\olref[fol][com][dls]{thm:downward-ls}). प्रथम, हे प्रमेय ज्या
तर्कशास्त्रांच्या सामान्य वर्गाला लागू होते, त्याचे अधिक व्यापक वर्णन
आपल्याला आवश्यक आहे. आपण स्वतःला \emph{संबंधात्मक} भाषांपुरते मर्यादित
ठेवू, म्हणजे ज्यांत केवळ !!{predicate}s आणि व्यक्ती सूचित करणारी स्थिर
चिन्हे असतात, पण !!{function}s नसतात अशा भाषा.

\end{document}
